\subsection{Proof of the relative-criterion comparison}
We prove Theorem~\ref{thm:criteria}, using the carrier calculus of
Proposition~5.7. Suppose the overlap restrictions in $\mathrm D$
are disjoint. Condition $\mathrm C$ makes $P,Q$ compatible.
Any common stratum would refine both parents and hence their
intersection, so it would belong to both restrictions. This is
impossible, proving $\mathrm C\Rightarrow\mathrm D$.
Taking $K=0$ gives $\mathrm D\Rightarrow\mathrm D_0$.
For distinct compatible elementary $x,y$ and
$p\rho x,q\rho y$, at least one of the two restrictions to
$x\cap y$ is zero. Indeed, if $x,y$ are incomparable, both
are zero; if $x<y$, the internal point of $y$ restricts to zero,
and the other ordered case is symmetric. Thus $\mathrm D$
gives $p\parallel q$, proving $\mathrm L3$.

The $\mathrm L2$ test extends from elementary objects to figures
by the stratum compatibility criterion and
$\ell(F)=\bigcup_{x\in F}\ell(x)$. Refinement decreases these
shadows, so $\mathrm L2$ implies $\mathrm D_0$.
Now assume $\mathrm L2$ and $\mathrm L3$ and take
$p\in\ell(P),q\in\ell(Q)$ under disjoint overlap restrictions.
Let $u=c_F(p),v=c_G(q)$. If $u\ne v$, these are distinct
compatible strata, and $p\rho u,q\rho v$ give disjunction by
$\mathrm L3$. If $u=v$, that carrier lies in $K$, so $p,q$
refine the respective overlap restrictions. Their disjunction
follows from $\mathrm L2$ there. Applying $\mathrm L2$ to all
such pairs proves $\mathrm D$.

For $\mathrm L3+\mathrm{CL}\Rightarrow\mathrm C$, suppose the
overlap restrictions are compatible. It suffices to show that
$x\in P,y\in Q$ are compatible. Put $H=x\cap y$ and take
$p\in\ell(x)\setminus\ell(H)$ and
$q\in\ell(y)\setminus\ell(H)$. Their carriers
$r=c_x(p),s=c_y(q)$ are different, since equality would put
$p$ below $H$. By carrier composition let
\[
 u=c_F(r)=c_F(p),\qquad v=c_G(s)=c_G(q).
\]
If $u\ne v$, they are distinct compatible strata of the parents,
and $\mathrm L3$ gives $p\parallel q$.
If $u=v$, that stratum lies in $K$, so $r\in P|K$ and
$s\in Q|K$. They are therefore compatible, and their distinctness
allows $\mathrm L3$ to give the same conclusion. The exterior
condition $\mathrm{CL}$ now proves $x\bowtie y$; compatibility
by strata proves $P\bowtie Q$.

Under $\mathrm D_0$, disjoint parents have disjoint point shadows.
Conversely, suppose all pairs from $\ell(x),\ell(y)$ are disjoint.
If $x\cap y$ contained a stratum, $\mathrm L1$ would supply a
lieu below it, disjoint from itself, a contradiction. Thus the
intersection is zero. Either $\mathrm{CL}$ or
$\mathrm{CL}_{\rm p}$ now gives compatibility and hence
disjointness. This proves the two remaining arrows to
$\mathrm L2$.

Under $\mathrm L2$ the two absolute antecedents are equivalent.
The $\mathrm{CD2}$ antecedent, tested on singleton lieux outside
$\ell(K)$, implies the exterior-shadow antecedent. Conversely,
if $P|K=0$, no lieu refining $P$ can refine $K$, by the
refinement-infimum property of restriction. The same holds for
$Q$. The exterior-shadow antecedent and figure-level $\mathrm L2$
then give $P\parallel Q$. The added polyhedral qualification is
unchanged. If an alternative polyhedral convention omits the
empty case, that case is already automatic under $\mathrm L2$:
its shadow test directly gives disjointness, hence compatibility.

Equal overlap restrictions are compatible, so $\mathrm C$ gives
the matching condition. For two matching refinements their union
is a figure, refines the union of their parents, and restricts to
the original pieces. Any other gluing is the union of its
restrictions and is equal to this one. Induction, or the
pairwise compatibility criterion together with restriction over
unions, proves finite gluing. Conversely the two-parent cover
gives the matching condition.

For an arbitrary cover, consider the increasing family of glued
refinements over its finite subcovers. Each is a refinement of
the full parent $F$. A stratum of $F$ belongs to some cover
member by primality. Once that member is included, matching
forces the restriction of every later finite gluing to equal
the prescribed local restriction. The stated closure hypothesis
therefore admits the union. It has exactly the required
restrictions and is unique. In a full clique atlas the union is
admitted directly, since every pair of its members is contained
in one of the finite gluings. The restricted-atlas example below
shows why this admission is a separate hypothesis.

\subsection{What the weakened stratum conditions miss}
The next examples concern explicitly stated order and disjunction
conditions. No ambiguous handwritten mark is assigned a definitive
reading by these calculations.

\begin{counterexample}[Least carriers without internal transitivity]
\label{ex:nontransitive}
On $\{p,q,z,u,t\}$, take incidence $z<t$ and refinement generated by
\[
 p\ll z,\quad p\ll u,\quad q\ll u,\quad u\ll t,
 \quad z\ll t,
\]
including identity pairs. Every principal deployment has least
carriers, and every open-lieu set is nonempty. With the two lieux
$p,q$ disjoint, the different strata in every principal deployment
have disjoint open-lieu sets. Nevertheless the induced internal
relation is not transitive, and disjunction defined by the open
lieux is not hereditary under it.
\end{counterexample}
\begin{proof}
All principal ideals are singletons except $[t]=\{z,t\}$.
Into $[t]$, the carriers of $p,z$ are $z$, and those of $q,u,t$
are $t$; every other carrier is forced by a singleton target.
The refinement-minimal objects are $p,q$, and
\[
 U_p=U_z=\{p\},\qquad U_q=U_t=\{q\},\qquad U_u=\{p,q\}.
\]
Thus all the stated point conditions hold, including the sole
nontrivial principal pair $z,t$. But
$p\rho u\rho t$ and $p\not\rho t$. The induced relation has
$p\delta t$ but not $p\delta u$, contradicting heredity along
$u\rho t$. This example is outside the strong-system hypotheses.
\end{proof}

\begin{counterexample}[Comparable pairs are insufficient]
\label{ex:weak-principal}
Partial incidence and internal orders, unique factorization,
exchange and hereditary symmetric irreflexive disjunction do not
force disjunction of every distinct pair in a principal deployment
if the stated disjunction requirement tests only comparable pairs.
The failure persists with nonempty open-lieu sets.
\end{counterexample}
\begin{proof}
First use $y,z<t$, internal identity, and only the edges $yt,zt$.
Internal factorization and exchange are identities; heredity is
automatic. Every comparable distinct pair is disjoint, but
$y,z<t$ are not. If instead the weakened clause requires comparable
pairs to be nondisjoint, the chain $a<t$ with internal identity and
empty disjunction satisfies it and again violates principal
disjunction.

For nonempty open lieux use $\{a,b,c,t\}$, incidence $a,b<t$,
and internal refinement $c\rho t$. The lieux are $a,b,c$,
with edges $ac,bc$ only. Their open sets are singletons, and
$U_t=\{c\}$. The disjunction induced by these open sets has
edges $ac,bc,at,bt$. Unique factors into $t$ are respectively
$a,b,t,t$ for the sources $a,b,c,t$; the only proper internal
arrow has singleton source, so exchange holds. Heredity follows
from the identical open sets of $c,t$. Both comparable pairs
are disjoint, but $a,b$ under $t$ are not. The missing condition
is precisely the all-pairs principal condition.
\end{proof}

\subsection{Four failures inside strong systems}
All four models below use the full family of incidence-closed
primitive cliques. Their internal relations are as stated or
computed from their carriers, so the general strong construction
supplies the prime core and its relative calculus.

\begin{counterexample}[Locality and extension need their hypotheses]
\label{ex:four-strong}
Strong systems alone do not imply $\mathrm L2$ or $\mathrm{CL}$,
even when all open-lieu sets are nonempty. The conditions
$\mathrm L2,\mathrm L3$ and witness absorption together do not
imply either absolute compatibility criterion. Subdivision
extension can also fail in a full strong atlas.
\end{counterexample}
\begin{proof}
For the first assertion take discrete incidence on $p,q,x,y$,
refinement generated by $p\ll x,q\ll y$, and sole primitive edge
$pq$. Principal deployments are singletons, and heredity only
tests the minimal edge endpoints, so this is strong. Both upper
open-lieu sets are nonempty, $U_x=\{p\},U_y=\{q\}$. Their points
are disjoint while $x,y$ are incompatible. Thus $\mathrm L2$
and $\mathrm{CL}$ fail at empty incidence intersection. Any
restricted intersection-class version allowing the empty
intersection also fails here.

A second direct failure of $\mathrm L2$ uses incidence and
refinement equal to the two chains $a<x,b<y$, with primitive
edges $ax,by,ab$. Internal refinement is identity; both principal
pairs are disjoint. The full lieu shadows of $x,y$ are
$\{a\},\{b\}$, which are disjoint, whereas their principal
figures are incompatible. The upper open-lieu sets are empty;
this second example is not asserted under $\mathrm L1$.

For failure of absolute compatibility take incidence and
refinement $a<x,a<y$, with primitive edges $ax,ay$.
Again internal refinement is identity and the system is strong.
Every nonempty figure contains $a$, so no two nonempty figures
are disjoint, and every full lieu shadow is $\{a\}$. This proves
$\mathrm L2$. Distinct compatible strata are $a$ with $x$ or
$y$; their upper open-lieu sets are empty, so $\mathrm L3$ holds.
All nonempty figures have the same cosupport $\{0\}$. Hence
witness absorption is vacuous for a nonzero refinement, and for
a zero refinement the parent is a witness. However the common
incidence figure of $x,y$ is $a$, and both exterior shadows are
empty. The absolute locus tests hold but $x,y$ are incompatible,
refuting $\mathrm{CL}$ and $\mathrm{CL}_{\rm p}$. Every nonzero
refinement of either parent restricts nontrivially to $a$, so the
outside-refinement tests hold as well, refuting $\mathrm{CD2}$
and $\mathrm{CD2}_{\rm p}$.

Finally let incidence and refinement be $a<b<c$, with isolated
$d$, and primitive edges $ab,ac,bc,ad,bd$, omitting $cd$.
Internal refinement is identity, and every principal ideal is
a clique. The complete list of figures is
\[
 0,\quad a,\quad d,\quad a+d,\quad a+b,\quad a+b+d,
 \quad a+b+c.
\]
Here sums denote unions of deployed strata. The figures disjoint
from $a$ and from $a+b$ are exactly $0,d$, so $a$ subdivides the
principal figure of $b$. Only $0$ is disjoint from $a+b+c$;
none of its proper refinements has that cosupport. Thus its only
subdivision is itself, and restriction to $b$ gives $a+b$, not
$a$. The prescribed subdivision of the subfigure does not extend.
This model has empty upper open-lieu sets, so it does not contradict
Theorem~\ref{thm:abstract-extension}.
\end{proof}

\subsection{Gluing does not give relative disjointness}
\begin{counterexample}[All covers can glue while $\mathrm C,\mathrm D$ fail]
\label{ex:gluing-relative}
There is a prime carrier atelier with $\mathrm L1,\mathrm L2$,
$\mathrm{CL}$, $\mathrm D_0$, witness absorption and unique gluing
over every set-indexed incidence cover, but without
$\mathrm C,\mathrm D$ or $\mathrm L3$.
\end{counterexample}
\begin{proof}
Use the four figures $0,\{x\},\{q\},\{x,y\}$.
Incidence is inclusion, and refinement adds
$\{q\}\ll\{x,y\}$. The three nonzero figures are completely
join-prime: a bounded family joining to one of them contains
that figure. Compatibility is equality together with the pair
of elementary strata $x,y$. All bounded joins and generation
are immediate. The least carriers are forced, with $q$ carried
by $y$ and $x$ by $x$ in the principal figure of $y$.
Compatible refinement is incidence. The sole proper internal
arrow is $q\rho y$, so internal transitivity holds. This verifies
all six prime-carrier conditions directly.

The lieux are $x,q$, with $U_y=\{q\}$, so $\mathrm L1$ holds.
No two nonempty figures are disjoint. The full shadows are
$\{x\},\{q\},\{x,q\}$; every two contain a nondisjoint point
pair, proving $\mathrm L2$. The condition $\mathrm D_0$ is
vacuous except when a parent is zero. For the only incompatible
elementary pairs, $x,q$ have nondisjoint exterior points and
$q,y$ have $q$ on both sides; thus their absolute test is false,
and $\mathrm{CL}$ holds. Every nonempty figure has cosupport
$\{0\}$, which proves witness absorption as in the preceding
example.

Every incidence cover of a nonzero figure contains that figure
itself by complete primality. Matching local data are therefore
determined by its chosen refinement, which is the unique gluing.
This holds for arbitrary covers, with repetitions allowed.
But for parents $x\leq y$, take refinements $x\ll x$ and
$q\ll y$. On their overlap $x$, the restrictions are $x$ and
$0$, which are disjoint and compatible. The refinements $x,q$
are neither. Thus $\mathrm C$ and $\mathrm D$ fail.
The compatible distinct strata $x,y$ have nondisjoint open lieux
$x,q$, so $\mathrm L3$ also fails. This is a prime-core example;
it is not claimed to be a strong canonical system.
\end{proof}

\subsection{The missing infinite atlas closure}
\begin{counterexample}[Finite gluing with an omitted stationary union]
\label{ex:restricted-infinite}
A restricted atlas can satisfy the prime core,
$\mathrm L1,\mathrm L2,\mathrm L3,\mathrm{CL}$ and finite gluing,
yet omit a directed union whose restrictions to every stratum
are eventually constant. It can consequently fail arbitrary-cover
gluing.
\end{counterexample}
\begin{proof}
Let $M=\{a_n,x_n:n\in\mathbb N\}$ have discrete incidence,
refinement $a_n\ll x_n$, and primitive disjunction exactly between
different indices. Each principal deployment is a singleton;
heredity preserves the index, so this is strong. The lieux are
$a_n$, and both the full and open lieu shadows of either member
of the $n$th pair are $\{a_n\}$. Distinguishing equal and unequal
indices verifies all four stated local conditions.

A full figure chooses at most one member of each pair. Admit
only figures with finitely many $a_n$, allowing arbitrarily
many $x_n$. This family contains all principal figures and zero,
is closed under subfigures, and is closed under every union
bounded by one of its admitted members. Finite compatible unions
are admitted too. Hence the strong restricted-atlas construction
and Theorem~\ref{thm:criteria} give the prime core and finite gluing.

The figure $F=\{x_n:n\in\mathbb N\}$ is admitted. The increasing
refinements $G_k=\{a_0,\ldots,a_k\}$ restrict to a fixed stratum
$x_n$ as $\{a_n\}$ once $k\geq n$. Their union has infinitely
many $a_n$ and is not admitted; no admitted incidence bound can
contain it. For a direct gluing failure cover $F$ by its singleton
strata and choose $\{a_n\}$ over $x_n$. These data match on their
empty overlaps, but any gluing would contain every $a_n$.
Thus no such gluing exists.
\end{proof}
